\exercisegroup

\begin{enumerate}
\def\labelenumi{\arabic{enumi})}
\item
  设 E 是向量空间，A 是 E 的一个对称凸子集。设 T、T' 是 E 上两个局部凸拓扑，U、U' 是 T、T' 在 E 上定义的一致结构。为使 U' 在 A 上诱导的一致结构细于 U 诱导的一致结构，其充要条件是：T 在 A 上诱导的拓扑的每个 0 的邻域，都是 T' 在 A 上诱导的拓扑的 0 的邻域。
\item
  \begin{enumerate}
  \def\labelenumii{\alph{enumii})}
  \tightlist
  \item
    给出 \(\mathbf{R}^2\) 中一个非紧闭集的例子，其凸包不是闭的。b) 证明：在 \(\mathbf{R}^n\) 中，紧集的凸包是紧的（参见 II, 第 66 页, 习题 9, a)）。
  \end{enumerate}
\end{enumerate}

¶ 3) 设 I 是 R 的紧区间 {[}0, 1{]}，F 是定义在 I 中的连续实值函数所成的向量空间 \(\mathcal{C}(\mathrm{I},\mathbf{R})\)。设 E 是乘积空间 \(\mathbf{R}^{\mathbf{F}}\)；对一切 \(a\in\mathrm{I}\)，设 \(\varepsilon_{a}\) 是 E 中使 \(\varepsilon_{a}(f)=f(a)\)（对一切 \(f\in\mathrm{F}\)）成立的元素。

\begin{enumerate}
\def\labelenumi{\alph{enumi})}
\item
  证明：当 \(x\) 在 I 中变化时，集 \(\mathbf{K}\)（由诸 \(\varepsilon_x\) 组成的）在 E 中是紧的。
\item
  设 \(\lambda\) 是 \(E\) 中使 \(\lambda(f) = \int_0^1 f(t)dt\)（对一切 \(f\in F\)）成立的元素（Lebesgue 测度）。
\end{enumerate}

证明：在 E 中，\(\lambda\) 属于 K 的凸包的闭包，但不属于该凸包（参见 FVR, II, 第 7 页, prop. 5）。

\begin{enumerate}
\def\labelenumi{\arabic{enumi})}
\setcounter{enumi}{3}
\tightlist
\item
  采用 II, 第 72 页, 习题 1 的记号，再假设空间 \(\mathbf{E}\) 是局部凸的。
\end{enumerate}

\begin{enumerate}
\def\labelenumi{\alph{enumi})}
\item
  存在一个正连续线性型 \(g\)（\(\mathbf{E}\) 中的），它延拓 \(f\)，其充要条件是：\(f\) 在 \(M \cap (\mathbf{W} + \mathbf{P})\) 中下有界，对至少一个邻域 \(\mathbf{W}\)（\(\mathbf{E}\) 中 0 的）而言。
\end{enumerate}

\begin{enumerate}
\def\labelenumi{\alph{enumi})}
\setcounter{enumi}{1}
\tightlist
\item
  给定一点 \(x \in \mathbf{E}\)，存在一个正连续线性型 \(g\)（\(\mathbf{E}\) 中的），使 \(g(x) = 1\)，其充要条件是 \(-x \notin \overline{\mathbf{P}}\)。
\end{enumerate}

\begin{enumerate}
\def\labelenumi{\arabic{enumi})}
\setcounter{enumi}{4}
\tightlist
\item
  \begin{enumerate}
  \def\labelenumii{\alph{enumii})}
  \tightlist
  \item
    设 E 是一个无穷维赋范空间，T 是它的拓扑。证明：E 上存在一个赋范空间拓扑 \(T'\)，它严格细于 T，又存在一个赋范空间拓扑 \(T''\)，它严格粗于 T（利用 E 的一个写成 \((a_{\alpha,n})\) 形式的基来定义这些拓扑的 0 的邻域，其中 \(\alpha\) 取遍一个无穷指标集 A，n 取遍整数 \(\geqslant 0\) 之集，且对 E 上给定的范数有 \(\|a_{\alpha,n}\| = 1\)）。
  \end{enumerate}
\end{enumerate}

\begin{enumerate}
\def\labelenumi{\alph{enumi})}
\setcounter{enumi}{1}
\item
  设 \(p\) 是定义拓扑 \(\mathcal{T}'\) 的范数。证明：若 \(E\) 对拓扑 \(\mathcal{T}\) 完备，则 \(p\) 在 \(E\) 中对拓扑 \(\mathcal{T}\) 不可能是下半连续的（利用 Baire 定理，参见 III, 第 25 页, 推论）。由此推出：凸集 \(A\)（由关系 \(p(x) < 1\) 定义的）对 \(\mathcal{T}\) 不含任何内点，尽管它的所有点都是代数内点。
\item
  由 \(b\) ) 推出：若 \(E\) 对拓扑 \(\mathcal{T}\) 完备，则 \(E\) 中存在对 \(\mathcal{T}\) 不是闭的凸集，它与每个有限维线性流形的交对 \(\mathcal{T}\) 都是闭的（参见 II, 第 67 页, 习题 12）。
\end{enumerate}

\begin{enumerate}
\def\labelenumi{\arabic{enumi})}
\setcounter{enumi}{5}
\tightlist
\item
  设 \(\mathbf{E}\) 是带其最细局部凸拓扑的向量空间。
\end{enumerate}

\begin{enumerate}
\def\labelenumi{\alph{enumi})}
\item
  证明：E 的每个向量子空间都是闭的，且若 M、N 是 E 中互为向量补的两个子空间，则 E 是 M 与 N 的拓扑直和。若 \((e_{t})_{t\in I}\) 是 E 的一个基，则 E 是诸子空间 \(Re_{t}\) 的拓扑直和。
\item
  设 \(\mathrm{F}\) 是一个局部凸空间，其拓扑也是最细局部凸拓扑。证明：\(\mathrm{E}\) 到 \(\mathrm{F}\) 中的每个线性映射都是严格态射。
\end{enumerate}

\begin{enumerate}
\def\labelenumi{\arabic{enumi})}
\setcounter{enumi}{6}
\item
  \begin{enumerate}
  \def\labelenumii{\alph{enumii})}
  \tightlist
  \item
    设 A 是拓扑向量空间 E 中至少含一个内点的凸集。证明：A 的代数内点所成之集等同于 A 的内部（参见习题 5, b)）。b) 证明：在赋范空间 \(E = l^{1}(\mathbf{N})\) 中，II, 第 66 页, 习题 11, b) 定义的凸锥 P 生成 E，但不含任何代数内点。
  \end{enumerate}
\item
  设 E 是带可数基与最细局部凸拓扑的向量空间。证明：若 E 中的集 A 与每个有限维向量子空间的交在 E 中是闭的，则 A 在 E 中是闭的（参见习题 5, c)）。
\item
  设 E 与 F 是两个向量空间，各带其最细局部凸拓扑。
\end{enumerate}

\begin{enumerate}
\def\labelenumi{\alph{enumi})}
\item
  证明：若 E 与 F 各有可数基，则 \(E \times F\) 到局部凸空间 G 中的每个双线性映射都是连续的（利用 Du Bois-Reymond 定理，FVR, V, 第 53 页, 习题 8）。
\item
  若空间 E、F 之一有一个基数等于连续统之基数的基，证明：\(E \times F\) 上存在一个不连续的双线性型。（化归到 \(E = R^{(N)}\)、\(F = R^{N}\) 的情形，于是 F 可等同于 \(E^{*}\)，且 \(E \times F\) 上的双线性型与 \(E^{*}\) 到自身的线性映射一一对应；然后考察 \(E^{*}\) 的恒等映射，并注意在 \(R^{N}\) 中，乘积拓扑下的紧集不可能是吸收的。）
\end{enumerate}

\begin{enumerate}
\def\labelenumi{\arabic{enumi})}
\setcounter{enumi}{9}
\item
  设 \((\mathrm{E}_n)\) 是局部凸空间的一个无穷序列，E 是族 \((\mathrm{E}_n)\) 的拓扑直和。证明：E 的拓扑等同于 I, 第 24 页, 习题 14 定义的拓扑 \(\mathcal{T}_0\)。
\item
  设 \(\mathbf{I}\) 是一个无穷不可数集。在向量空间 \(\mathrm{E} = \mathbf{R}^{(I)}\) 上，证明：最细局部凸拓扑不同于拓扑 \(\mathcal{T}_0\)（定义见 \(\mathbf{I}\), 第 24 页, 习题 14）；为此证明：由点 \(x = (\xi_{i})_{i\in I} \in \mathrm{E}\)（满足 \(|\sum_{i\in I} \xi_i| < 1\)）所成之集在 \(\mathcal{T}\) 中是开的，但在 \(\mathcal{T}_0\) 中不是。
\item
  设 \(\mathbf{E}\) 是带可数基 \((e_n)\) 的向量空间。设 \(\mathbf{V}\) 是诸 \(e_n\) 所成之集的均衡凸包，\(\mathbf{W}\) 是由点
\end{enumerate}

\[
a _ {n} = e _ {n} + (n - 1) e _ {1} \quad (n \geqslant 1).
\]

设 \(T_{1}\)（相应地 \(T_{2}\)）是 E 上以诸 \(\lambda V\)（相应地 \(\lambda W\)）（\(\lambda > 0\)）为 0 的基本邻域系的局部凸拓扑。证明：\(T_{1}\) 与 \(T_{2}\) 都是 Hausdorff 的，但 \(T_{1}\) 与 \(T_{2}\) 在 E 上局部凸拓扑之集中的下确界不是 Hausdorff 的（参见 II, 第 80 页, 习题 26）。

\begin{enumerate}
\def\labelenumi{\arabic{enumi})}
\setcounter{enumi}{12}
\item
  在 II, §6 的假设下，证明：E 对归纳极限为诸 \(T_{n}\) 的拓扑 T 是完备的，其充要条件是：对每个整数 n 以及 \(E_{n}\) 上每个对 T 诱导的拓扑而言的 Cauchy 滤子 F，存在 \(p \geqslant n\)，使 F 在 \(E_{p}\) 中对拓扑 \(T_{p}\) 收敛。
\item
  设 E 是局部凸空间递增序列 \(E_{n}\)（II, 第 33 页）的严格归纳极限。证明：E 的拓扑是 E 上与向量空间结构相容的拓扑（无论是否局部凸）中在 \(E_{n}\) 上诱导出粗于给定拓扑 \(T_{n}\) 的拓扑者之中的最细者。（设 \(V_{0}\) 是这样一个拓扑 T 的 0 的邻域，\((\mathrm{V}_{n})_{n\geqslant0}\) 是 T 的 0 的邻域序列，使 \(V_{n+1} + V_{n+1} \subset V_{n}\) 对一切 \(n \geqslant 0\) 成立；对一切 \(n \geqslant 1\)，在 \(E_{n}\) 中考察 0 的一个凸邻域 \(W_{n}\)（含于 \(E_{n} \cap V_{n}\)），并取诸 \(W_{n}\) 在 E 中之并的凸包。）
\item
  设 I 是一个无穷不可数集。设 \(\mathfrak{F}(\mathrm{I})\) 是 I 的有限子集所成的族，E 是直和空间 \(\mathbf{R}^{(I)}\)。对每个 \(J \in \mathfrak{F}(\mathrm{I})\)，设 \(F_{J}\) 是 E 的子空间 \(\mathbf{R}^{\mathrm{J}}\)（指标属于 J 的诸因子之积，赋以乘积拓扑）；设 \(g_{J}\) 是 \(F_{J}\) 到 E 中的典范单射。证明：存在一个拓扑 \(\mathcal{T}_0\)（\(E\) 上的），它与 \(E\) 的向量空间结构相容，使诸 \(g_{J}\) 连续，且严格细于最细局部凸拓扑 \(\mathcal{T}\)（它使诸 \(g_{J}\) 连续）。（注意：集 \(V\)——由点 \(x = (\xi_{t})_{t\in I}\)（\(E\) 中满足 \(\sum_{i\in I}|\xi_i|^{1 / 2}\leqslant 1\) 的）组成——是某个与向量
\end{enumerate}

空间结构相容的拓扑的 0 的邻域，且该集不含任何吸收的对称凸集。）

\begin{enumerate}
\def\labelenumi{\arabic{enumi})}
\setcounter{enumi}{15}
\tightlist
\item
  \begin{enumerate}
  \def\labelenumii{\alph{enumii})}
  \tightlist
  \item
    设 E 是带可数基的向量空间。证明：E 上最细局部凸拓扑是 E 上在 E 的每个有限维子空间上诱导出典范拓扑的拓扑（无论是否与 E 的向量空间结构相容）中的最细者。
  \end{enumerate}
\end{enumerate}

\begin{enumerate}
\def\labelenumi{\alph{enumi})}
\setcounter{enumi}{1}
\tightlist
\item
  设 \(E_{0}\) 是一个无穷维 Banach 空间。设 E 是 \(E_{0}\) 与 \(\mathbf{R}^{(\mathbf{N})}\) 的直和向量空间，\(E_{p}\) 是 E 的子空间，即 \(E_{0}\) 与 \(R^{p}\) 的直和（等同于 \(\mathbf{R}^{(\mathbf{N})}\) 前 p 个因子之积；我们给 \(E_{p}\) 赋以其诸因子拓扑的乘积拓扑，使 \(E_{p}\) 的拓扑由 \(E_{p+1}\) 的拓扑诱导）。证明：在 E 上，诸 \(E_{p}\) 的拓扑的归纳极限拓扑并不是 E 上在每个 \(E_{p}\) 上诱导出粗于 \(E_{p}\) 之拓扑的拓扑（无论是否与 E 的向量空间结构相容）中的最细者。可如下进行：
\end{enumerate}

\(\alpha)\) 设 \(q\) 是 \(\mathbf{E}_0\) 上一个范数，它定义一个严格粗于 \(\mathbf{E}_0\) 之拓扑的拓扑（II, 第 74 页, 习题 5）。对每个 \(\varepsilon > 0\)，定义一个映射 \(f_{\varepsilon}\)（\(\mathbf{E}_0\) 到 \(\mathbf{R}_+\) 中的）如下：\(f_{\varepsilon}(x) = \sup(q(x), \varepsilon - \| x \|)\)。证明：\(f_{\varepsilon}\) 连续且 \(>0\)（在 \(\mathbf{E}_0\) 中），并且

\[
\inf_{\| x \| = \varepsilon} f_{\varepsilon}(x) = 0 .
\]

\(\beta)\) 设 \(U\) 是 \(E\) 中由诸 \((x, (t_n))\)（满足 \(t_n < f_{1/n}(x)\)，对一切 \(n\)）所成的子集。证明：\(U \cap E_p\) 在 \(E_p\) 中对一切 \(p\) 是开的。

\(\gamma)\) 证明：若 \(\mathrm{V} \subset \mathrm{U}\) 是吸收凸集，则 \(\mathrm{V} \cap \mathrm{E}_0\) 不能包含 \(\mathrm{E}_0\) 中任何以 0 为中心的球。

\begin{enumerate}
\def\labelenumi{\arabic{enumi})}
\setcounter{enumi}{16}
\tightlist
\item
  对交换群 G（写成加法）的子集 A 以及每个 \(n > 0\)，把形如 \(\sum_{i=1}^{n} x_i\)（其中 \(x_i \in \mathbf{A}\) 对一切 \(i\) 成立）的元素所成之集记作 \({}^{n}+\mathbf{A}\)。我们说集 A
\end{enumerate}

是凸的，如果对每个整数 n > 0，关系 \(nx \in {}^{n}+A\) 蕴含 \(x \in A\)。

\begin{enumerate}
\def\labelenumi{\alph{enumi})}
\item
  证明：若交换拓扑群 G（写成加法）同构于一个局部凸向量空间（带诱导拓扑）的加法群的一个子群，则 G 中存在对称凸的 0 的基本邻域系。
\item
  反之，设 G 是一个 Hausdorff 交换拓扑群（写成加法），其中存在对称凸的 0 的基本邻域系 B。证明：G 是无挠的，因而可（在代数上）看作 Q 上向量空间的加法群的一个子群（A, II, § 7.10, prop. 26 的推论 1）。对每个集 V ∈ B，设 \(\tilde{V}\) 是形如 rx 的元素所成之集，其中 x ∈ V，而 r 取遍满足 \(0 \leqslant r \leqslant 1\) 的有理数；证明：\(\tilde{V}\) 是对称且凸的（按上面定义的意义）。进一步推出：若 G 没有异于 G 自身的开子群，则诸集 \(\tilde{V}\) 构成 Q 上 E 的与向量空间结构相容的拓扑（Q 赋以其通常拓扑）的 0 的基本邻域系；由此断言，此时 G 同构于一个 Hausdorff 局部凸空间的加法群的一个子群。
\item
  设 G 是按字典序排序的群 \(R \times R\)（A, VI, 第 7 页）；考察 G 上与其群结构相容的 Hausdorff 拓扑 \(\mathcal{T}_{0}(G)\)（GT, IV, § 1, 习题 1）。证明：对这个拓扑存在对称凸的 0 的基本邻域系，但 G 不同构于 R 上任何 Hausdorff 拓扑向量空间的加法群的任何子群。
\end{enumerate}

